\chapter{Testiranje i evaluacija}

\section{Automatizovano testiranje backend-a}

Backend je testiran automatizovanim, tabelarno-vođenim (\emph{table-driven}) Go testovima, standardnim pristupom u Go ekosistemu gde se više varijacija istog scenarija testira kroz jednu test funkciju sa nizom ulaznih/izlaznih parova. Tabela 9.1 prikazuje broj test funkcija po paketu, na dan pisanja rada; svi navedeni testovi prolaze (\texttt{go test ./...} vraća \texttt{ok} za svaki paket koji sadrži testove).

\setlength{\tabcolsep}{4pt}
\renewcommand{\arraystretch}{1.3}

\begin{longtable}{@{}|>{\small\raggedright\arraybackslash}p{0.26\textwidth}|>{\small\centering\arraybackslash}p{0.11\textwidth}|>{\small\raggedright\arraybackslash}p{0.51\textwidth}|@{}}
\caption{Broj automatizovanih testova po Go paketu}\label{tab:9-1}\\
\hline
\rowcolor{black!12}
\textbf{Paket} & \textbf{Broj testova} & \textbf{Fokus testiranja} \\
\hline
\endhead
\texttt{internal/\allowbreak algorithm} & 15 & Isključivanje nedopustivih ivica/čvorova, poklapanje Dijkstra/A*, penal za skretanje, ponašanje nad realnim OSM podacima \\
\hline
\texttt{internal/scoring} & 6 & Funkcija bodovanja rute, preferenca korisnika, popust za omiljenu stanicu \\
\hline
\texttt{internal/explain} & 8 & Detekcija odstupanja rute, identifikacija vezujuće dimenzije, geometrijska divergencija \\
\hline
\texttt{internal/reststop} & 9 & Pretraga najbližeg odmarališta u koridoru rute, hazmat preferenca, brend/omiljena lokacija \\
\hline
\texttt{internal/auth} & 5 & Izdavanje/parsiranje JWT tokena, provera lozinke, verzija tokena \\
\hline
\texttt{internal/ws} & 4 & Relej pozicije pretplatnicima, replay poslednje pozicije novom pretplatniku \\
\hline
\texttt{internal/httpapi} & 4 & HTTP handleri (integracioni testovi nad REST rutama) \\
\hline
\texttt{internal/geocode} & 6 & Nominatim proxy, throttling \\
\hline
\textbf{Ukupno} & \textbf{57} & \\
\hline
\end{longtable}

\setlength{\tabcolsep}{6pt}

Paket \texttt{internal/algorithm} koristi dve vrste test podataka: (a) male, ručno konstruisane sintetičke grafove (npr. "dijamant" od četiri čvora sa kontrolisanim ograničenjima), pogodne za preciznu proveru jedne izolovane osobine algoritma (npr. "grana sa \texttt{maxheight=3.5} se isključuje za vozilo visine 4.0m"), i (b) stvaran OSM ekstrakt koridora Beograd–Novi Sad (poglavlje 9.2), koji dodatno potvrđuje da implementacija korektno parsira i koristi \emph{stvarne} OSM tagove, ne samo tagove koje je autor sam smislio za potrebe testa.

\section{Evaluacija sopstvenog algoritma nad realnim podacima}

Sopstvena Dijkstra/A* implementacija (poglavlje 6.4) je evaluirana nad stvarnim OSM ekstraktom koridora Beograd–Novi Sad (major-roads-only ekstrakt, \newline \texttt{testdata/beograd-novisad-corridor.osm}), učitanim istim \texttt{LoadOSMXML} mehanizmom koji koristi i produkcioni kod. Ekstrakt se učitava u graf od \textbf{38 265 čvorova i 50 705 usmerenih ivica} — dva reda veličine manje od kompletnog nacionalnog Valhalla grafa (898 337 čvorova, poglavlje 5.2), u skladu sa namerom da ovaj modul bude ograničen, transparentan alat za evaluaciju, ne zamena za Valhallu.

\textbf{Plauzibilnost rute.} Za par tačaka Beograd (44.80, 20.40) – Novi Sad (45.25, 19.85), Dijkstra pretraga (vozilo visine 4.0m, mase 40 000 kg) pronalazi put od \textbf{81,0 km kroz 1899 čvorova}. Za poređenje, Valhalla, nad kompletnim nacionalnim grafom sa punim skupom puteva (ne samo magistralnim), za isti par tačaka vraća rute u rasponu od 84,8 do 92,7 km (u zavisnosti od izabrane alternative), zabeleženo tokom rada na Valhalla integraciji (poglavlje 6.1). Rezultat sopstvenog modula je unutar plauzibilnog opsega — blago kraći, što je i očekivano jer ograničeni ekstrakt sadrži samo magistralne puteve i ne modeluje sve raskrsnice/petlje koje Valhalla uzima u obzir sa kompletnim grafom, a start/cilj pretrage su najbliži čvorovi grafa vazdušnoj liniji do date koordinate, ne tačno iste ulazno-izlazne rampe koje bi Valhalla izabrala.

\textbf{Isključivanje na osnovu stvarnog OSM tag-a.} Nad istim ekstraktom pronađena je stvarna deonica sa tagom \texttt{maxheight=4.3} u centralnoj zoni Beograda. Za vozilo visine 4.0m, Dijkstra pronalazi put od 4036m preko te deonice; za vozilo visine 4.5m (koje premašuje 4,3m ograničenje), pretraga korektno ne pronalazi nijedan put između istih dveju tačaka — u ovom ograničenom, samo-magistralnom ekstraktu ne postoji alternativna deonica koja bi vozilo moglo koristiti umesto isključene, pa je ishod potpuno odsustvo puta, a ne duži zaobilazni put. Ovaj test potvrđuje da mehanizam isključivanja radi korektno na \textbf{stvarnim}, ne samo ručno konstruisanim OSM podacima.

\textbf{Poklapanje Dijkstra i A*\emph{.} Za isti par (Beograd, Novi Sad) i profil, Dijkstra i A*} implementacija vraćaju identičnu cenu puta, što je i teorijski očekivano za dopustivu heuristiku (poglavlje 2.3) — A* ne pronalazi drugačiji (a kamoli lošiji) put, samo ga pronalazi pregledavajući manji broj čvorova. Vremenska analiza (Tabela 9.2) potvrđuje da je pretraga, nad grafom ove veličine, brza u odnosu na jednokratno parsiranje ulaznog XML fajla — 11,3 MB OSM XML-a.

\setlength{\tabcolsep}{4pt}
\renewcommand{\arraystretch}{1.3}

\begin{longtable}{@{}|>{\small\raggedright\arraybackslash}p{0.62\textwidth}|>{\small\centering\arraybackslash}p{0.20\textwidth}|@{}}
\caption{Vreme izvršavanja nad koridorom Beograd–Novi Sad (38 265 čvorova, 50 705 ivica)}\label{tab:9-2}\\
\hline
\rowcolor{black!12}
\textbf{Korak} & \textbf{Vreme} \\
\hline
\endhead
Učitavanje i parsiranje OSM XML ekstrakta (\texttt{LoadOSMXML}) & ≈ 405 ms \\
\hline
Dijkstra pretraga (Beograd → Novi Sad) & ≈ 19,3 ms \\
\hline
A* pretraga (Beograd → Novi Sad) & ≈ 17,1 ms \\
\hline
\end{longtable}

\setlength{\tabcolsep}{6pt}

Rezultat u Tabeli 9.2 potvrđuje očekivanje iz poglavlja 2.3 i 6.4: sama pretraga, jednom kada je graf učitan u memoriju, izvršava se za desetine milisekundi i dominantno je ograničena parsiranjem ulaznih podataka, a ne samim algoritmom pretrage — što je i razlog zbog kojeg je u produkcionom putu ovog sistema izgradnja i održavanje grafa u celosti prepušteno Valhalla-i (poglavlje 5.2, 6.1), koja tu istu vrstu posla radi jednom, unapred, za ceo nacionalni graf, umesto pri svakom zahtevu.

\section{Uticaj prilagođene funkcije rizika — studija slučaja}

Doprinos sopstvenog sloja rangiranja (poglavlje 6.2) u odnosu na "sirov" izbor prve Valhalla alternative ilustrovan je konkretnim slučajem otkrivenim tokom testiranja, na ruti u okolini mesta Radalj. Rana verzija formule bodovanja, koja nije sadržala vremenski član (\texttt{timeTerm}), birala je kao "najbolju" alternativu rutu koja je bila \textbf{47\% duža i 30\% sporija} od druge ponuđene alternative, isključivo zato što je imala povoljniji odnos magistralne i lokalne deonice puta (\texttt{highwayTerm}). Sa stanovišta stvarnog transportnog troška (vreme, gorivo, habanje vozila), taj izbor je bio pogrešan — kraća/brža alternativa je bila objektivno bolji izbor uprkos manje povoljnom \texttt{highwayRatio}.

Dodavanjem \texttt{timeTerm}-a, koji direktno kažnjava kandidata srazmerno njegovom relativnom kašnjenju u odnosu na najbrži kandidat u istom skupu alternativa (poglavlje 6.2, Listing 6.3), formula je počela korektno da favorizuje bržu/kraću rutu u ovom slučaju, bez uklanjanja postojećih članova formule (broj manevara, udeo magistralne deonice, potrošnja, oštri manevri) koji ostaju relevantni u slučajevima kada razlika u trajanju nije toliko izražena. Ovaj slučaj je dobar primer opšte lekcije o dizajnu heurističkih funkcija cene sa više članova: dodavanje novog signala treba motivisati \textbf{konkretnim, uočenim lošim ishodom} postojeće formule, a ne unapred pretpostavljenom kompletnošću — isti princip po kome je i ova formula, i funkcija cene u poglavlju 6.4, iz istog razloga ostala eksplicitno obeležena kao nekalibrisana, prva heuristička procena, a ne konačan, empirijski potvrđen model.

\section{Ograničenja identifikovana testiranjem}

Testiranje je otkrilo i granice sistema koje nisu greške u uobičajenom smislu, već posledice svesno odabranog obima rada:

\begin{itemize}
\item Sopstveni Dijkstra/A* modul (poglavlje 6.4), primenjen na isti par koordinata za koji je poznat slučaj Novi Banovci (poglavlje 6.3), \textbf{ne reprodukuje} taj slučaj — u ograničenom, samo-magistralnom ekstraktu korišćenom za ovaj modul, u okolini Novih Banovaca ne postoji \texttt{maxheight} tag na way nivou. Najverovatnije objašnjenje, dokumentovano tokom razvoja, je da je stvarno ograničenje u tom slučaju zavedeno kao tačkasta (node) prepreka, ne kao way tag, i/ili da geometrijski najkraći put u ograničenom grafu jednostavno ne prolazi kroz tačno tu deonicu koju bi Valhalla, sa svojim potpuno drugačijim, kompletnim grafom i sopstvenom heuristikom, izabrala. Ovo je pošteno prijavljen negativan nalaz, ne prikriven — dodatno naglašava razliku u nameni dva sloja rutiranja u ovom sistemu: Valhalla (sa Explain mehanizmom, poglavlje 6.3) je taj koji stvarno rešava slučaj Novi Banovci u produkciji; sopstveni modul (6.4) demonstrira princip mehanizma isključivanja nad drugim, ali i dalje stvarnim, primerom istog tipa OSM ograničenja (maxheight=4,3 u centralnom Beogradu, poglavlje 9.2).
\item Penal za skretanje u sopstvenoj pretrazi (6.4) je, kako je i dokumentovano u samom kodu, teorijsko pojednostavljenje koje u retkim slučajevima ne garantuje globalno optimalan put po ukupnom penalu za skretanje.
\item Funkcija rizika (6.2) i funkcija cene sopstvene pretrage (6.4) koriste bazne koeficijente koji su prva heuristička procena, ne kalibrisani prema stvarnim podacima o potrošnji goriva, habanju vozila ili statistici saobraćajnih nezgoda — razlika u odnosu na produkcioni sistem bi bila prikupljanje takvih podataka i kalibracija koeficijenata regresijom ili sličnim metodom, što izlazi iz obima ovog rada.
\end{itemize}

\clearpage
